\section{三大序列处理原语对比}

在本课程中，我们已经学习了三种不同的神经网络原语来处理序列/结构化数据：

\begin{figure}[H]
\centering
\includegraphics[width=0.95\textwidth]{slides-images/slide-043.jpg}
\caption{RNN vs CNN vs Self Attention 的对比：各自擅长的数据类型、并行性和扩展性\protect\footnotemark}
\end{figure}
\footnotetext{Slides 第44页。}

\begin{table}[H]
\centering
\begin{tabular}{lccc}
\toprule
 & \textbf{RNN} & \textbf{CNN} & \textbf{Self Attention} \\
\midrule
\textbf{操作对象} & 1D 有序序列 & 多维网格 & 无序向量集合 \\
\textbf{信息交互范围} & 理论无限 & 局部（核大小） & 全局（一层即可） \\
\textbf{序列方向并行性} & 不可 & 可以 & 可以 \\
\textbf{计算复杂度} & $O(n)$ & $O(n)$ & $O(n^2)$ \\
\textbf{长距离依赖} & 困难（梯度消失） & 需要多层 & 一层直达 \\
\textbf{可扩展性} & 差 & 中 & 好 \\
\bottomrule
\end{tabular}
\caption{三种序列处理原语的对比}
\end{table}

\begin{warningbox}{Self Attention 的 $O(n^2)$ 代价}
Self Attention 的计算和内存复杂度都是 $O(n^2)$（$n$ 为序列长度）：
\begin{itemize}
    \item 需要计算所有 $n \times n$ 对的相似度
    \item 需要存储 $n \times n$ 的注意力矩阵
\end{itemize}
当 $n = 100{,}000$ 甚至 $n = 1{,}000{,}000$ 时，这会变得非常昂贵。但 Justin Johnson 指出，\textbf{更多计算并不一定是坏事}——更多的计算意味着网络有更强的``思考''能力。解决方案是：买更多 GPU。
\end{warningbox}

\begin{importantbox}{Attention is All You Need}
在三种原语中，Self Attention 凭借以下优势胜出：
\begin{enumerate}
    \item \textbf{全局信息交互}：一层 Self Attention 就能让所有位置互相``看到''
    \item \textbf{高度并行}：核心计算是矩阵乘法，完美适配 GPU
    \item \textbf{可扩展性}：通过增加硬件即可处理更大的模型和更长的序列
\end{enumerate}
这正是 2017 年那篇标志性论文的标题所传达的信息：``Attention is All You Need''。
\end{importantbox}

\subsection{本章小结}

RNN 擅长序列但不可并行，CNN 擅长网格但感受野有限，Self Attention 擅长全局交互但代价是 $O(n^2)$ 复杂度。在现代深度学习的规模下，Self Attention 的可扩展性使其成为首选。

%% ========== Transformer 架构 ==========

